\chapter*{Danh mục các từ viết tắt}
\addcontentsline{toc}{chapter}{Danh mục các từ viết tắt}

{\small
\begin{longtable}{r l >{\raggedright\arraybackslash}p{10.6cm}}
\toprule
\textbf{STT} & \textbf{Viết tắt} & \textbf{Tên đầy đủ} \\
\midrule
\endfirsthead
\toprule
\textbf{STT} & \textbf{Viết tắt} & \textbf{Tên đầy đủ} \\
\midrule
\endhead
\bottomrule
\endfoot
1 & AI & Artificial Intelligence -- Trí tuệ nhân tạo \\
2 & API & Application Programming Interface -- Giao diện lập trình ứng dụng \\
3 & AUC & Area Under the ROC Curve -- Diện tích dưới đường cong ROC \\
4 & CanFG & Cancelable Face Generator -- Bộ sinh khuôn mặt có thể hủy \\
5 & CFB & Cancelable Face Biometrics -- Sinh trắc học khuôn mặt có thể hủy \\
6 & DL & Deep Learning -- Học sâu \\
7 & DoS & Denial of Service -- Tấn công từ chối dịch vụ \\
8 & e4e & Encoder for Editing -- Bộ mã hóa phục vụ chỉnh sửa \\
9 & EER & Equal Error Rate -- Tỷ lệ lỗi cân bằng \\
10 & FAR & False Accept Rate -- Tỷ lệ chấp nhận sai \\
11 & FaRL & Facial Representation Learning -- Học biểu diễn khuôn mặt đa phương thức \\
12 & FID & Fréchet Inception Distance -- Khoảng cách Fréchet Inception \\
13 & FMR & False Match Rate -- Tỷ lệ chấp nhận nhầm \\
14 & FNMR & False Non-Match Rate -- Tỷ lệ từ chối nhầm \\
15 & FR & Face Recognition -- Nhận dạng khuôn mặt \\
16 & FRR & False Reject Rate -- Tỷ lệ từ chối sai \\
17 & GAN & Generative Adversarial Network -- Mạng đối sinh \\
18 & ISO/IEC & International Organization for Standardization / International Electrotechnical Commission -- Tổ chức Tiêu chuẩn hóa Quốc tế / Ủy ban Kỹ thuật Điện Quốc tế \\
19 & MLP & Multi-Layer Perceptron -- Mạng nơ-ron đa tầng \\
20 & MTCNN & Multi-task Cascaded Convolutional Networks -- Mạng tích chập phân tầng đa nhiệm \\
21 & OTB-Morph & One-Time Biometrics via Morphing -- Sinh trắc học dùng một lần bằng lai ghép ảnh \\
22 & PSR & Protection Success Rate -- Tỷ lệ bảo vệ thành công \\
23 & ROC & Receiver Operating Characteristic -- Đường cong đặc trưng hoạt động bộ thu \\
24 & SOTA & State-of-the-Art -- Tiên tiến nhất hiện nay \\
25 & StyleGAN & Style-based Generative Adversarial Network -- Mạng đối sinh dựa trên phong cách \\
26 & TPS & Thin-Plate Spline -- Nội suy bản mỏng \\
27 & ViT & Vision Transformer -- Mô hình Transformer cho thị giác máy tính \\
\end{longtable}
\addtocounter{table}{-1}
}

\cleardoublepage
\chapter*{Danh mục các thuật ngữ}
\addcontentsline{toc}{chapter}{Danh mục các thuật ngữ}

{\small
\begin{longtable}{r >{\raggedright\arraybackslash}p{5.2cm} >{\raggedright\arraybackslash}p{8.0cm}}
\toprule
\textbf{STT} & \textbf{Tiếng Anh} & \textbf{Tiếng Việt} \\
\midrule
\endfirsthead
\toprule
\textbf{STT} & \textbf{Tiếng Anh} & \textbf{Tiếng Việt} \\
\midrule
\endhead
\bottomrule
\endfoot
1 & Anonymity & Tính ẩn danh \\
2 & Artifact & Biến dạng giả / Nhiễu nhân tạo \\
3 & Attack Tree & Cây tấn công \\
4 & Attribute Preservation & Bảo toàn thuộc tính (phi danh tính) \\
5 & Augmentation & Tăng cường dữ liệu / Tạo biến thể \\
6 & Biometric System Security & Bảo mật hệ thống sinh trắc học \\
7 & Black-box & Kịch bản hộp đen \\
8 & Cancelable Biometrics & Sinh trắc học có thể hủy \\
9 & Confidence Score & Điểm tin cậy (tương đồng) \\
10 & Control Grid & Lưới điểm neo / Lưới điều khiển \\
11 & Cross-matching & Đối chiếu chéo (liên kết danh tính) \\
12 & Disentanglement & Tách biệt thuộc tính (trong không gian tiềm ẩn) \\
13 & Distortion Scheme & Lược đồ biến dạng \\
14 & Enrollment & Pha đăng ký mẫu sinh trắc \\
15 & Face Morphing & Lai ghép khuôn mặt \\
16 & Feature-space & Không gian đặc trưng \\
17 & GAN Inversion & Đảo ngược GAN (chiếu ảnh vào không gian tiềm ẩn) \\
18 & Grid Morphing & Biến dạng lưới hình học \\
19 & Hill-Climbing & Thuật toán leo đồi \\
20 & Identity Preservation & Bảo toàn danh tính \\
21 & Image Morphing & Biến dạng ảnh / Lai ghép ảnh \\
22 & Image-space & Không gian ảnh (thô) \\
23 & Irreversibility & Tính không thể đảo ngược \\
24 & Key Mapping & Ánh xạ khóa \\
25 & Latent Code & Mã tiềm ẩn \\
26 & Latent Space & Không gian tiềm ẩn \\
27 & Loss Function & Hàm mất mát \\
28 & Na\"ive Mixing & Phép trộn thô / Trộn khởi tạo \\
29 & Optimization Loop & Vòng lặp tối ưu hóa \\
30 & Performance Preservation & Bảo toàn hiệu năng (nhận dạng) \\
31 & Pre-trained & Tiền huấn luyện \\
32 & Protected Face & Khuôn mặt được bảo vệ \\
33 & Random Seed & Hạt giống ngẫu nhiên \\
34 & Revocability & Tính khả hủy (thu hồi và cấp lại) \\
35 & Revocable Biometrics & Sinh trắc học có thể thu hồi / hủy \\
36 & Skip-connection & Liên kết tắt \\
37 & Style Mixing & Hòa trộn phong cách \\
38 & Test-time Optimization & Tối ưu hóa ở pha suy luận \\
39 & Thin-Plate Spline & Nội suy bản mỏng \\
40 & Undistort & Giải biến dạng ngược \\
41 & Unlinkability & Tính không thể liên kết \\
42 & Warping & Biến dạng hình học ảnh \\
\end{longtable}
\addtocounter{table}{-1}
}

\cleardoublepage
\chapter*{Danh mục các ký hiệu}
\addcontentsline{toc}{chapter}{Danh mục các ký hiệu}

{\small
\begin{longtable}{r l >{\raggedright\arraybackslash}p{8.8cm}}
\toprule
\textbf{STT} & \textbf{Ký hiệu} & \textbf{Ý nghĩa và vai trò} \\
\midrule
\endfirsthead
\toprule
\textbf{STT} & \textbf{Ký hiệu} & \textbf{Ý nghĩa và vai trò} \\
\midrule
\endhead
\bottomrule
\endfoot
1 & $\mathcal{X}$ & Không gian ảnh khuôn mặt. \\
2 & $\mathcal{K}$ & Không gian khóa có thể hủy. \\
3 & $\mathcal{W}^+$ & Không gian tiềm ẩn trung gian phân tầng của StyleGAN2. \\
4 & $\mathbf{x}_r,\; \mathbf{x}_p \in \mathcal{X}$ & Ảnh khuôn mặt thật đầu vào và ảnh được bảo vệ đầu ra. \\
5 & $\mathbf{k} \in \mathcal{K}$ & Khóa bí mật dùng làm hạt giống ngẫu nhiên. \\
6 & $\mathbf{z}_r,\; \mathbf{z}_f,\; \mathbf{z}_p \in \mathcal{W}^+$ & Mã tiềm ẩn của ảnh thật, mã sinh từ khóa và mã sau khi hòa trộn. \\
7 & $\mathbf{F} : \mathcal{X} \times \mathcal{X} \rightarrow \{0, 1\}$ & Bộ so khớp khuôn mặt ($1$: cùng người $\mathbf{x}_1 \equiv \mathbf{x}_2$; $0$: khác người $\mathbf{x}_1 \not\equiv \mathbf{x}_2$). \\
8 & $\mathbf{G} : \mathcal{X} \times \mathcal{K} \rightarrow \mathcal{X}$ & Bộ sinh khuôn mặt có thể hủy cần xây dựng ($\mathbf{x}_p = \mathbf{G}(\mathbf{x}_r, \mathbf{k})$). \\
9 & $\mathbf{S} : \mathcal{K} \rightarrow \mathcal{W}^+$ & Hàm ánh xạ khóa $\mathbf{k}$ sang mã tiềm ẩn giả lập $\mathbf{z}_f$. \\
10 & $\mathbf{E} : \mathcal{X} \rightarrow \mathcal{W}^+$ & Bộ mã hóa đảo ngược \textit{e4e} chiếu ảnh thật $\mathbf{x}_r$ sang mã tiềm ẩn $\mathbf{z}_r$. \\
11 & $\mathbf{M} : \mathcal{W}^+ \times \mathcal{W}^+ \rightarrow \mathcal{W}^+$ & Hàm hòa trộn $\mathbf{z}_r$ và $\mathbf{z}_f$ thành mã tiềm ẩn trộn $\mathbf{z}_p$. \\
12 & $\mathbf{D} : \mathcal{W}^+ \rightarrow \mathcal{X}$ & Bộ sinh StyleGAN2 tiền huấn luyện giải mã mã tiềm ẩn về ảnh khuôn mặt. \\
13 & $\mathbf{I},\; \mathbf{A}$ & Bộ trích xuất đặc trưng danh tính (ArcFace) và phi danh tính (FaRL). \\
14 & $\mathcal{L}_{\text{anon}}, \mathcal{L}_{\text{idp}}, \mathcal{L}_{\text{attr}}$ & Các hàm mất mát ẩn danh, bảo toàn danh tính và bảo toàn thuộc tính. \\
15 & $I(x, y),\; I'$ & Ảnh khuôn mặt gốc và ảnh mẫu sống khi xác thực. \\
16 & $\mathbf{p}_{i,j},\; \Delta \mathbf{p}_{i,j}(\mathbf{k})$ & Tọa độ điểm nút lưới điều khiển và độ dịch chuyển theo khóa $\mathbf{k}$. \\
17 & $\mathcal{T}_{\mathbf{k}},\; \mathcal{T}_{\mathbf{k}}^{-1}$ & Phép biến đổi tọa độ liên tục theo khóa $\mathbf{k}$ và ánh xạ ngược. \\
18 & $I_{\mathbf{k}},\; I_{\text{key}},\; I_p$ & Ảnh sau biến dạng lưới, ảnh tham chiếu theo khóa và ảnh sau pha trộn lai ghép. \\
19 & $\alpha \in (0, 1)$ & Hệ số pha trộn tuyến tính trong lai ghép khuôn mặt. \\
20 & $\mathcal{E}(\cdot),\; s(\cdot, \cdot),\; \tau$ & Bộ trích xuất đặc trưng chuẩn, hàm đo độ tương đồng và ngưỡng quyết định xác thực. \\
\end{longtable}
\addtocounter{table}{-1}
}
